\section{Limitations and Conclusion}

\subsection{Limitations}

The claim boundary in Section~2 and Table~\ref{tab:limitations} applies throughout. The evidence uses synthetic agent tools and scenarios, receipt-based correlation, one reference gateway, local-harness performance, and one frozen model/runtime selection. The post-denial substitution study evaluates predeclared branches in a bounded reference runtime, not deployment prevalence. The Campaign Correlation frontier establishes selectivity only for its declared deterministic receiver-local strategy, not a general escalation policy or deployment effectiveness. The full external-warning study consumes preserved detector outputs for \ExternalWarningFullSourceCases{} source cases: \ExternalWarningFullWarnings{} produce admitted exact evidence and \ExternalWarningFullAttackNoWarning{} attack-bearing cases produce no warning. It does not estimate detector quality, warning prevalence, arbitrary warning-source behavior, or production effectiveness. The protocol lifecycle lacks evidence for key rotation, revocation propagation, heterogeneous federation deployment, and cross-host transport delay.

The validation-value study establishes a precise, narrow result. Its protected population contains \MatchingHarmfulActions{} accepted binding-matched harmful actions. C1-class and C2 both block those actions, while C2 avoids the \COneClassBenignConstrained{} benign constraints observed under C1-class. Across the full cohort, C1-class blocks \COneClassHarmfulBlocked{} of \ValidationEpisodes{} harmful actions and C2 blocks \CTwoHarmfulBlocked{}; C2 rejects mismatched, expired, replayed, or uncorrelated warnings rather than applying a broad class guard. The hostile-peer, correlation, and mutation studies are deterministic synthetic sweeps. They establish behavior in the examined reference path rather than a field prevalence estimate.

The retained amendment history makes the RA-03 eligibility correction visible. Table~\ref{tab:ra03-eligibility} reports the before-and-after rule outcomes. The correction affects the benign availability cohort; the strict deterministic causal comparison excludes RA-03 and remains separately rooted.

\subsection{Conclusion}

TCOP specifies an evidence-to-authorization boundary at which cross-domain runtime evidence becomes an input to a receiver's local authorization decision for a pending agent action. The main result is conditional on timeliness: when accepted evidence reaches the receiver inside an action's containment window, it can alter locally governed enforcement and reduce later harmful agent-action opportunities. TCOP specifies authenticated evidence exchange, receiver acceptance, and the authority boundary. It does not prescribe Exact Binding, Campaign Correlation, monitoring, or a standing guard. In every tested single-environment reference-path and replay condition, gateway blocks cite receiver-local policy and decision records.

The validation-value study shows why the full evidence binding matters. On the accepted binding-matched harmful population, validated matching preserves the class guard's containment while avoiding its measured collateral benign constraints. The post-denial substitution and Campaign Correlation frontier studies show alternative receiver strategies supported by the same boundary: Exact Binding prioritizes precision, while the evaluated Campaign Correlation strategy covers the declared related substitutions and forwards the measured outside-campaign same-risk actions. The external-warning study separates source coverage from authorization: after declared normalization and local receipt construction, receiver policies act on accepted evidence, while no-warning cases remain outside conditional TCOP coverage. The artifact preserves the source roots, exclusions, sealed study plans, normalized data, and regenerated figures and tables so reviewers can inspect each current claim against its evidence.
